\section{Evaluation}
\label{sec:evaluation}

\todolist{tim}{
\item show Sudoku constraints 
\item explain constraints
\item add multiplayer events with constraints (to demonstrate how a
  new codepath does not break the system and ``just works''}
  
  
This section evaluates how well constraint applications with Babelsberg/S facilitate reactive systems.

We argued that logic puzzles, like Sudoku puzzles, are good examples for interactive constraint applications. Therefore, the evaluation uses Sudoku as example application. In \autoref{lst:sudoku-constraints}, you will find all necessary constraints to solve a Sudoku puzzle.

\begin{lstlisting}[caption=All Constraints of a Sudoku Puzzle, label=lst:sudoku-constraints]
sudoku addConstraintsForAllGivenNumbers.

[sudoku allSatisfy: [:cell | cell between: 1 and: 9]] 
		alwaysSolveWith: solver.
		
(1 to: sudoku rowCount) do: [:index |
		[(sudoku atRow: index) allDifferent &
		(sudoku atColumn: index) allDifferent &
		(sudoku atBox: index) allDifferent] 
				alwaysSolveWith: solver].
\end{lstlisting}

The second one, line 3 and 4, was already explained in \autoref{sec:ocp_squeak}. The first line adds equality constraints for all given numbers, so that the constraint solver cannot change them. In general, they are called stay constraints, which are currently not supported in Z3, but will be in future versions. The last constraint, line 6 to 10, describes that all numbers in each row, column and box, are different. The keyword for that is \texttt{allDifferent}, which is a new collection method in Babelsberg/S. Together, the second and third constraint say that each row, column and box contains all numbers from 1 to 9.